\documentclass[11pt]{article}
\usepackage[T1]{fontenc}
\usepackage[utf8]{inputenc}
\usepackage{fourier}
\usepackage[scaled=0.875]{helvet}
\renewcommand{\ttdefault}{lmtt}
\usepackage{amsmath,amssymb}
\usepackage{makeidx}
\makeindex
\usepackage{fancybox}
\usepackage[normalem]{ulem}
\usepackage{enumitem}
\usepackage{pifont}
\newcommand{\textding}[1]{\text{\ding{#1}}}
\usepackage{lscape}
\usepackage{lastpage}
\usepackage{tabularx,booktabs}
\usepackage{scratch3}
\usepackage{textcomp}
\usepackage{enumitem}
\usepackage{multicol}
\usepackage{multirow}
\usepackage{mdframed}
\definecolor{scrmovedddd}{HTML}{3373cc}%
\newcommand{\euro}{\eurologo{}}
%Tapuscrit : Denis Vergès
\usepackage{pst-eucl}
\usepackage{diagbox}% à mettre après pst-eucl
\usepackage{graphicx,pstricks,pst-plot,pst-grad,pst-node,pst-text,pstricks-add}
\usepackage{pstricks-add}
\usepackage{tikz}
\usetikzlibrary{arrows.meta, calc, angles, quotes}
\usepackage{pgfplots}
\pgfplotsset{compat=1.18}
\definecolor{gris}{gray}{0.85}
\usetikzlibrary{patterns,calc,decorations.pathmorphing}
\newcommand{\R}{\mathbb{R}}
\newcommand{\N}{\mathbb{N}}
\newcommand{\D}{\mathbb{D}}
\newcommand{\Z}{\mathbb{Z}}
\newcommand{\Q}{\mathbb{Q}}
\newcommand{\C}{\mathbb{C}}
\usepackage[left=3.5cm, right=3.5cm, top=2.2cm, bottom=2cm]{geometry}
\setlength\headheight{13.7pt}
\newcommand{\vect}[1]{\overrightarrow{\,\mathstrut#1\,}}
\renewcommand{\theenumi}{\textbf{\arabic{enumi}}}
\renewcommand{\labelenumi}{\textbf{\theenumi.}}
\renewcommand{\theenumii}{\textbf{\alph{enumii}}}
\renewcommand{\labelenumii}{\textbf{\theenumii.}}
\def\Oij{$\left(\text{O}~;~\vect{\imath},~\vect{\jmath}\right)$}
\def\Oijk{$\left(\text{O}~;~\vect{\imath},~\vect{\jmath},~\vect{k}\right)$}
\def\Ouv{$\left(\text{O}~;~\vect{u},~\vect{v}\right)$}
%\newcommand{\arc}[1]{\widehat{#1}}
\usepackage{fancyhdr}
\usepackage[french]{babel}
\usepackage[np]{numprint}
\usepackage{colortbl}
\usepackage[colorlinks=true,linkcolor=blue,citecolor=blue,urlcolor=blue]{hyperref}
\hypersetup{%
pdfauthor = {APMEP},
pdfsubject = {Brevet},
pdftitle = {Année 2026},
allbordercolors = white,
pdfstartview=FitH}
\begin{document}
\setlength\parindent{0mm}
\marginpar{\rotatebox{90}{\textbf{A. P{}. M. E. P{}.}}}
\lhead{\small L'année 2026}
\rhead{\small \textbf{A. P{}. M. E. P{}.}}
\rfoot{\small Amérique du Nord}
\lfoot{\small 3 juin 2026}
\pagestyle{fancy}
\thispagestyle{empty}
\begin{center} {\Large \textbf{\decofourleft~Diplôme national du brevet Amérique du Nord~\decofourright\\[7pt] 3 juin 2026}}
\end{center}

%\section*{Exercice 1 :\hfill 20 points}

%\begin{center}
%  {\large\textbf{DIPLÔME NATIONAL DU BREVET}}\\[4pt]
%  {\large\textbf{SESSION 2026}}\\[20pt]
%  \fbox{\parbox{0.75\textwidth}{\centering
%    \vspace{12pt}
%    {\LARGE\textbf{MATHÉMATIQUES}}\\[10pt]
%    {\large\textbf{Série générale}}\\[12pt]
%    \textbf{Durée : 2 h 00} \hfill \textbf{Coefficient : 2}\\[8pt]
%  }}
%\end{center}
%
%\vspace{12pt}
%\begin{center}
%  Dès que le sujet vous est remis, assurez-vous qu'il est complet.\\
%  Ce sujet comporte 8 pages numérotées de la page 1/8 à la page 8/8.
%\end{center}

%\vspace{12pt}
\begin{center}
\begin{tabular}{|p{0.60\textwidth}|p{0.25\textwidth}|}
\hline
\textbf{Partie 1 – Automatismes} \newline 20 min (calculatrice interdite) & \textbf{6 points} \\
\hline
\textbf{Partie 2 – Raisonnement et résolution de problèmes} \newline 1 h 40 (calculatrice autorisée) & \textbf{14 points} \\
\hline
\end{tabular}
\end{center}

\vspace{12pt}

\begin{center}
\textbf{À l'issue de la partie 1, les copies sont ramassées.}
\end{center}

\vspace{12pt}
L'usage de la calculatrice avec mode examen actif ou sans mémoire \og type collège \fg est \textbf{interdit} pour la partie 1 et \textbf{autorisé} pour la partie 2.

L'utilisation du dictionnaire est interdite.

\bigskip


\section*{Partie 1 – Automatismes – 6 points – 20 minutes}


\begin{mdframed}
\itshape
Pour chaque question, recopier sur la copie son numéro et la réponse correspondante.\\
Pour cette partie, aucune justification n'est demandée.\\
Pour les questions à choix multiple, une seule réponse est exacte.
\end{mdframed}

\vspace{10pt}

\paragraph{Question 1}
Calculer $A = \dfrac{2}{3} + \dfrac{3}{4}$.

\vspace{8pt}
\paragraph{Question 2}
Un article coûte 45\,\euro. Quel sera son prix après une réduction de 10\,\%\,?

\vspace{8pt}
\paragraph{Question 3}
Un professeur a dessiné à main levée le quadrilatère ci-dessous avec ses diagonales.\\
Que peut-on affirmer à propos de la nature de ce quadrilatère\,?\\
Recopier sur la copie la lettre de la bonne réponse.

\vspace{6pt}
%Tikz
%\begin{center}
%\begin{tikzpicture}[scale=1.1]
%  % Quadrilatère quelconque avec marques sur les côtés opposés
%  \coordinate (A) at (0,0);
%  \coordinate (B) at (3.2,0.3);
%  \coordinate (C) at (3.5,1.8);
%  \coordinate (D) at (0.4,1.5);
%  % Diagonales
%  \draw (A) -- (B) -- (C) -- (D) -- cycle;
%  \draw (A) -- (C);
%  \draw (B) -- (D);
%  % Marques égalité côtés AB et CD
%  \draw[thick] (1.55,0.13) -- (1.65,0.17);
%  \draw[thick] (1.85,1.65) -- (1.95,1.69);
%  % Marques égalité côtés BC et DA
%  \draw[thick] (3.4,1.05) -- (3.37,1.15);
%  \draw[thick] (0.15,0.78) -- (0.12,0.88);
%\end{tikzpicture}
%\end{center}
%Fin Tikz
%PsTricks
\begin{center}
\psset{unit=1.5cm}
\begin{pspicture}(4,2)
\uput[dl](0,0){A}\uput[dr](3.2,0.3){B}\uput[ur](3.5,1.8){C}\uput[ul](0.4,1.5){D}
\def\barre{\psline(-0.05,-0.1)(-0.05,0.1)\psline(0.05,-0.1)(0.05,0.1)}
\rput{30}(1,0.5){\barre}\rput{30}(2.6,1.35){\barre}\rput{-30}(1.1,1.2){\barre}\rput{-30}(2.4,0.65){\barre}
\pslineByHand(0,0)(0.4,1.5)(3.5,1.8)(3.2,0.3)(0,0)
\pslineByHand(0,0)(3.5,1.8)
\pslineByHand(0.4,1.5)(3.2,0.3)
\end{pspicture}
\end{center}

\vspace{6pt}
\begin{center}
\begin{tabular}{|p{0.19\textwidth}|p{0.19\textwidth}|p{0.19\textwidth}|p{0.25\textwidth}|}
\hline
\centering\textbf{Réponse A} & \centering\textbf{Réponse B} & \centering\textbf{Réponse C} & \centering\textbf{Réponse D}\tabularnewline
\hline
\centering C'est un losange & \centering C'est un rectangle & \centering C'est un carré &
\centering Ce n'est ni un losange, ni un rectangle\tabularnewline
\hline
\end{tabular}
\end{center}

\vspace{8pt}
\paragraph{Question 4}
Résoudre l'équation $5x - 15 = 20$.

\medskip

\paragraph{Question 5}
Dans le repère ci-contre, on a placé deux points A et B.

\medskip
\begin{minipage}{0.45\textwidth}
\begin{enumerate}[label=\textbf{\alph*.}]
\item Quelle est l'abscisse du point A\,?
\item Quelles sont les coordonnées du point B\,?
\end{enumerate}
\end{minipage}\hfill
\begin{minipage}{0.52\textwidth}
\begin{center}
%Tikz
%\begin{tikzpicture}[scale=0.65]
%\begin{axis}[
%axis lines=center,
%xmin=-5.5, xmax=3.5,
%ymin=-3.5, ymax=3.5,
%xtick={-5,-4,...,3},
%ytick={-3,-2,...,3},
%tick label style={font=\small},
%xlabel={},ylabel={},
%grid=both,
%grid style={line width=0.3pt, draw=gray!40},
%width=7.5cm, height=7.5cm,
%]
%\addplot[only marks, mark=x, mark size=3pt, thick] coordinates {(-2,2)};
%\node[above left, font=\small] at (axis cs:-2,2) {A};
%\addplot[only marks, mark=x, mark size=3pt, thick] coordinates {(-2,-1)};
%\node[below right, font=\small] at (axis cs:-2,-1) {B};
%\end{axis}
%\end{tikzpicture}
%fin Tikz
%PSTricks
\psset{unit=1cm,arrowsize=2pt 3}
\begin{pspicture}(-4.1,-3.1)(3.1,3.1)
\psaxes[linewidth=1.25pt,labelFontSize=\scriptstyle]{->}(0,0)(-4.1,-3.1)(3.1,3.1)
\psgrid[subgriddiv=1,gridlabels=0pt,gridwidth=0.15pt](-4.,-3.)(3.,3.)
\psdots(-2,2)(-2,-1)
\uput[l](-2,2){A}\uput[l](-2,-1){B}
\end{pspicture}
%Fin PSTricks
\end{center}
\end{minipage}

\vspace{10pt}
\paragraph{Question 6}
Voici une série de nombres : 8 ; 19 ; 12 ; 3 ; 12 ; 25 ; 3 ; 11 ; 1.\\
Déterminer la médiane de cette série.

\vspace{8pt}
\paragraph{Question 7}
On considère un triangle ABC rectangle en A tel que :
\begin{itemize}
\item BC $= 5$ cm
\item $\widehat{\text{ABC}} = 60\degres$.
\end{itemize}
Recopier sur la copie la formule qui permet d'obtenir la longueur AB.

\vspace{6pt}
\begin{minipage}{0.45\textwidth}
\begin{center}
\begin{tabular}{|c|c|}
\hline
$5 \times \sin(60)$ & $5 \times \cos(60)$ \\
\hline
$5 \div \sin(60)$ & $5 \div \cos(60)$ \\
\hline
\end{tabular}
\end{center}
\end{minipage}
\hfill
\begin{minipage}{0.50\textwidth}
\begin{center}
%Tikz
%\begin{tikzpicture}[scale=1.0]
%\coordinate (A) at (0,0);
%\coordinate (B) at (2.5,0);
%\coordinate (C) at (0,2.5);
%\draw[thick] (A) -- (B) -- (C) -- cycle;
%% Angle droit en A
%  \draw (0.25,0) -- (0.25,0.25) -- (0,0.25);
%  % Angle 60° en B
%  \node[below right] at (B) {$60°$};
%  % Labels
%  \node[below left] at (A) {A};
%  \node[below right, xshift=2pt] at (B) {};
%  \node[below] at (B) {B};
%  \node[above left] at (C) {C};
%  % mesure BC
%  \node[right] at ($(B)!0.5!(C)$) {$5$ cm};
%\end{tikzpicture}
%Fin Tikz
%PSTricks
\psset{unit=1cm,arrowsize=2pt 3}
\begin{pspicture}(7,4.5)
\pspolygon(2.6,0.8)(4.5,0.8)(2.6,4)%ABC
\uput[dl](2.6,0.8){A}\uput[dr](4.5,0.8){B}\uput[l](2.6,4){C}
\psarc(4.5,0.8){0.4cm}{120}{180}\psframe(2.6,0.8)(2.8,1)
\uput[ul](4.1,1){$60\degres$}\uput[ur](3.55,2.4){5 cm}
\rput(3.5,0.1){\small\emph{La figure n'est pas représentée en vraie grandeur}}
\end{pspicture}
%Fin PSTricks
\vspace{4pt}

\end{center}
\end{minipage}

\vspace{10pt}
\paragraph{Question 8}
Donner un diviseur de 387 autre que 1 et lui-même.

\vspace{20pt}
\begin{center}
  \textbf{Restitution de la copie du candidat à l'issue de la partie 1}
\end{center}

\newpage

%%%%%%%%%%%%%%%%%%
\section*{Partie 2 – Raisonnement et résolution de problèmes 14 points 1 h 40}


\begin{mdframed}
\itshape
Dans cette partie, toutes les réponses doivent être justifiées, sauf si une indication contraire est donnée.\\
La clarté et la précision des raisonnements ainsi que la rédaction sont évaluées sur 2 points.\\
Pour chaque question, si le travail n'est pas terminé, laisser tout de même une trace de la
recherche ; les essais et les démarches engagées, même non aboutis, seront pris en compte dans la notation.
\end{mdframed}

\vspace{10pt}

%-------------------------------------------------------------
\subsection*{Exercice 1 \normalfont(2,5 points)}
%-------------------------------------------------------------

La figure ci-contre n'est pas représentée en vraie grandeur.

\vspace{6pt}
\begin{minipage}[t]{0.48\textwidth}
\vspace{0pt}
Les points B, A et E sont alignés.\\
Les points C, A et D sont alignés.\\
Le triangle ABC est rectangle en B.
\begin{itemize}
\item DE $= 4{,}8$ cm
\item AD $= 7{,}3$ cm
\item AE $= 5{,}5$ cm
\item BC $= 7{,}2$ cm.
\end{itemize}
\end{minipage}
\hfill
\begin{minipage}[t]{0.48\textwidth}
\vspace{0pt}
\begin{center}
%début Tikz
%\begin{tikzpicture}[scale=0.9]
%% Coordonnées
%\coordinate (B) at (0,0);
%\coordinate (A) at (1.8,0);
%\coordinate (E) at (4.2,0);
%\coordinate (C) at (1.8,3.5);
%\coordinate (D) at (3.5,-2.0);
%% Triangle ABC
%\draw[thick] (B) -- (A) -- (C) -- cycle;
%% Angle droit en B
%\draw (0.25,0) -- (0.25,0.25) -- (0,0.25);
%% Extension vers E et D
%\draw[thick] (A) -- (E);
%\draw[thick] (A) -- (D);
%% Triangle AED
%\draw[thick] (E) -- (D);
%% Angle droit en E
%\draw ($(E)!0.3cm!(A)$) -- ++(0,-0.3) -- ($(E)+(0,-0.3)$);
%% Labels
%\node[below left] at (B) {B};
%\node[above] at (C) {C};
%\node[above right, xshift=2pt] at (A) {A};
%\node[right] at (E) {E};
%\node[below right] at (D) {D};
%\end{tikzpicture}
%Fin Tikz
%PSTricks
\psset{unit=1cm,arrowsize=2pt 3}
\begin{pspicture}(5.8,7)
\pspolygon(0.4,3)(5.4,3)(5.4,0.2)(0.4,6.2)%BEDC
\uput[ur](3,3){A}\uput[dl](0.4,3){B}\uput[ur](5.4,3){E}\uput[dr](5.4,0.2){D}
\uput[ul](0.4,6.2){C}
\psframe(0.4,3)(0.6,3.2)
\end{pspicture}
%Fin PSTricks
\end{center}
\end{minipage}

\vspace{10pt}

\begin{enumerate}
\item Montrer que le triangle AED est un triangle rectangle en E.
\item Calculer l'aire du triangle AED.
\item Pourquoi peut-on affirmer que les droites (BC) et (ED) sont parallèles ?
\item Calculer la valeur exacte de la longueur AB.
\item On admet que l'angle $\widehat{\text{ACB}}$ mesure environ $49\degres$. En déduire la mesure de l'angle $\widehat{\text{ADE}}$.
\end{enumerate}

\newpage

%-------------------------------------------------------------
\subsection*{Exercice 2 \normalfont(3,5 points)}
%-------------------------------------------------------------

On considère les fonctions $f$ et $g$ définies par :
$f(x) = (x-1)(x+3)$ \quad et \quad $g(x) = 2x+1$.

\begin{enumerate}
\item Calculer $f(-4)$.
\item Déterminer l'antécédent de 2 par la fonction $g$.
\item On utilise un tableur pour donner les images des nombres entiers de 0 à 8 par les fonctions $f$ et $g$.

\vspace{6pt}
\begin{center}
\begin{tabularx}{\linewidth}{|c|*{10}{>{\centering \arraybackslash}X|}}\hline
& \textbf{A} & \textbf{B} & \textbf{C} & \textbf{D} & \textbf{E} & \textbf{F} & \textbf{G} & \textbf{H} & \textbf{I} & \textbf{J} \\\hline
\textbf{1} & $x$ & 0 & 1 & 2 & 3 & 4 & 5 & 6 & 7 & 8 \\\hline
\textbf{2} & $f(x)$ & $-3$ & 0 & 5 & 12 & 21 & 32 & 45 & 60 & 77 \\\hline
\textbf{3} & $g(x)$ & 1 & 3 & 5 & 7 & 9 & 11 & 13 & 15 & 17 \\\hline
\end{tabularx}
\end{center}

\begin{enumerate}[label=\alph*)]
\item Quelle formule doit-on saisir en cellule B3 puis étirer vers la droite pour compléter la ligne 3\,?\\
\textit{Aucune justification n'est demandée.}
\item Par lecture du tableau ci-dessus, donner une solution de l'équation $f(x) = g(x)$.\\
\textit{Aucune justification n'est demandée.}
\end{enumerate}

\item On représente graphiquement chacune de ces fonctions.

\medskip

\begin{minipage}[t]{0.52\textwidth}
\begin{center}
%début Tikz
%\begin{tikzpicture}
%\begin{axis}[
%axis lines=center,
%xmin=-4.5, xmax=3.5,
%ymin=-5.5, ymax=5.5,
%xtick={-4,-3,...,3},
%ytick={-5,-4,...,5},
%tick label style={font=\small},
%grid=both,
%grid style={line width=0.3pt, draw=gray!40},
%width=7.5cm, height=8.5cm,
%]
%% C1 = courbe parabolique f(x)=(x-1)(x+3)
%\addplot[thick, domain=-4.2:3.2, samples=80] {(x-1)*(x+3)};
%\node[above left, font=\small] at (axis cs:-3.8,4.5) {$\mathcal{C}_1$};
%% C2 = droite g(x)=2x+1
%\addplot[thick, domain=-4.2:3.2, samples=2] {2*x+1};
%\node[below right, font=\small] at (axis cs:2,4) {$\mathcal{C}_2$};
%\end{axis}
%\end{tikzpicture}
%Fin Tikz
%début PSTricks
\psset{unit=1cm,arrowsize=2pt 3}
\begin{pspicture*}(-4.5,-5.5)(3.5,5.5)
\psaxes[linewidth=1.25pt,labelFontSize=\scriptstyle]{->}(0,0)(-4.5,-5.5)(3.5,5.5)
\psplot[plotpoints=2000,linewidth=1.25pt,linecolor=red]{-4}{2.5}{x 1 sub x 3 add mul}
\uput[r](-3.5,3.5){\red $\mathcal{C}_1$}
\psplot[plotpoints=200,linewidth=1.25pt]{-4}{2.5}{x 2 mul 1 add}\uput[r](-2.7,-4.5){ $\mathcal{C}_2$}
\end{pspicture*}
%fin PSTricks
\end{center}
\end{minipage}
\hfill
\begin{minipage}[t]{0.44\textwidth}
\vspace{20pt}
\begin{enumerate}[label=\alph*)]
\item Associer à chacune des fonctions $f$ et $g$ sa représentation graphique.\\
\textit{Aucune justification n'est demandée.}
\item Par lecture graphique, déterminer les deux solutions de l'équation $f(x) = g(x)$.\\
\textit{Aucune justification n'est demandée.}
\end{enumerate}
\end{minipage}

\item Lola affirme que les solutions de l'équation $f(x) = g(x)$ sont les mêmes que les solutions de l'équation $x^2 - 4 = 0$. A-t-elle raison\,? Justifier.
\end{enumerate}

\newpage

%%%%%%%%%
\subsection*{Exercice 3 \normalfont(4 points)}


Dans cet exercice, les deux parties sont indépendantes.

Une entreprise développe une intelligence artificielle (IA) capable de reconnaître des objets sur des images.

\subsubsection*{Partie A}

On entraîne l'IA à partir d'une base de données de \np{50000} images réparties en 4 catégories : \og Objets du quotidien \fg, \og Animaux \fg, \og Véhicules \fg, \og Autres \fg.

L'intelligence artificielle est testée pour mesurer sa précision et son efficacité. Les images sont réparties comme suit :

\vspace{6pt}
\begin{center}
\begin{tabular}{|>{\columncolor{gris}}c|l|c|}\hline
\rowcolor{gris}
& \textbf{Type d'image} & \textbf{Nombre d'images} \\\hline
2 & Objets du quotidien & 28\,000 \\\hline
3 & Animaux & \np{12000} \\\hline
4 & Véhicules & \np{8000} \\\hline
5 & Autres & ? \\\hline
\end{tabular}
\end{center}

\medskip

\begin{enumerate}
\item Combien d'images appartiennent à la catégorie \og Autres \fg\,?
\item Sur l'ensemble des tests, l'intelligence artificielle reconnaît correctement 90\,\% des \og Objets du quotidien \fg.

Calculer le nombre d'images reconnues correctement dans cette catégorie.
\item L'intelligence artificielle reconnaît correctement \np{5600} images de la catégorie \og Véhicules \fg.

Quel pourcentage de réussite cela représente-t-il dans cette catégorie\,?
\item Une image est tirée au hasard dans la base de données.\\
Quelle est la probabilité que l'image tirée soit l'image d'un \og Objet du quotidien \fg\,?

On donnera le résultat sous la forme d'un nombre décimal.
\end{enumerate}

\bigskip

\subsubsection*{Partie B}

L'intelligence artificielle, très utilisée dans le monde entier, nécessite une quantité importante d'électricité. L'énergie consommée peut s'exprimer en wattheures (Wh).

En 2024, sa consommation annuelle est estimée à 82\,000 Gigawattheures (GWh).

En comparaison, un collège consomme en moyenne 200\,000 kilowattheures (kWh) par an.

\medskip

\begin{center}
\begin{tabularx}{0.35\linewidth}{|X|}\hline
\textbf{Rappels :}\\
-- 1\,kWh $= 10^3$\,Wh\\
-- 1\,GWh $= 10^9$\,Wh\\ \hline
\end{tabularx}
\end{center}

\medskip

\begin{enumerate}[resume]
\item Convertir la consommation de l'IA et d'un collège en Wh.

\emph{Exprimer ces résultats sous la forme d'une écriture scientifique.}
\item Combien de collèges pourrait-on alimenter pendant un an avec la consommation électrique de l'intelligence artificielle\,?
\item En France, il y a environ \np{7100} collèges. Dans cette question, on suppose que chaque collège a la même consommation d'énergie annuelle moyenne (\np{200000} kWh).

Pendant combien d'années environ pourrait-on alimenter tous les collèges français avec la consommation électrique annuelle de cette intelligence artificielle\,?
\end{enumerate}

%\newpage

%%%%%%%%%%%%%%
\subsection*{Exercice 4 \normalfont(2 points)}

Dans cet exercice, aucune justification n'est demandée.

Un élève souhaite réaliser une figure constituée de carrés et de triangles équilatéraux, à l'aide d'un logiciel de programmation. Pour cela, il crée les trois blocs ci-dessous :

\vspace{10pt}

\begin{center}
\begin{tabularx}{\linewidth}{|*{3}{X|}}\hline
\multicolumn{1}{|c|}{\textbf{Bloc 1}} &
\multicolumn{1}{|c|}{\textbf{Bloc 2}} &\multicolumn{1}{|c|}{\textbf{Bloc 3}}\\ \hline
\smallskip
\begin{scratch}[scale=0.86]
\initmoreblocks{définir \namemoreblocks{Initialisation}}
\blockmove{aller à x : \ovalnum{0} y : \ovalnum{0}}
\blockmove{s'orienter à ~\ovalnum{90}}
\blockpen{effacer tout}
\blockpen{stylo en position d'écriture}
\end{scratch}
&
\smallskip
\begin{scratch}[scale=0.86]
\initmoreblocks{définir \namemoreblocks{carré}}
\blockrepeat{répéter \ovalnum{\textbf{A}} fois}
{
\blockmove{avancer de \ovalnum{50} pas}
\blockmove{tourner \turnleft{} de \ovalnum{\textbf{B}} degrés}
}
\end{scratch}
&
\smallskip
\begin{scratch}[scale=0.86]
\initmoreblocks{définir \namemoreblocks{triangle équilatéral}}
\blockrepeat{répéter \ovalnum{\textbf{C}} fois}
{
\blockmove{avancer de \ovalnum{50} pas}
\blockmove{tourner \turnright{} de \ovalnum{\textbf{D}} degrés}
}
\end{scratch}
\\ \hline
\end{tabularx}
\end{center}

\medskip


L'instruction \og s'orienter à 90 \fg{} signifie que le lutin se dirige vers la droite.

\begin{enumerate}
\item Quelles sont les coordonnées du lutin après l'exécution du Bloc 1\,?
\item Dans les blocs 2 et 3, on a remplacé certaines valeurs par les lettres \textbf{A}, \textbf{B}, \textbf{C} et \textbf{D}.

Sur la copie, indiquer la lettre et sa valeur correspondante.
\item L'élève a construit trois figures avec les trois programmes ci-dessous.

Associer chaque figure au programme correspondant.
\end{enumerate}

\medskip

\begin{center}
\begin{tabularx}{\linewidth}{|*{3}{X|}}\hline
\multicolumn{1}{|c|}{\textbf{Programme 1}} &
\multicolumn{1}{|c|}{\textbf{Programme 2}} &\multicolumn{1}{|c|}{\textbf{Programme 3}}\\ \hline
\begin{scratch}[scale=0.9]
\blockinit{quand \greenflag est cliqué}
\blockmoreblocks{Initialisation}
\blockmoreblocks{triangle équilatéral}
\blockrepeat{répéter \ovalnum{\textbf{3}} fois}
{
\blockmoreblocks{carré}
\blockmove{avancer de \ovalnum{50} pas}
\blockmove{tourner \turnright{} de \ovalnum{\textbf{120}} degrés}
}
\end{scratch}
&
\begin{scratch}[scale=0.9]
\blockinit{quand \greenflag est cliqué}
\blockmoreblocks{Initialisation}
\blockmoreblocks{carré}
\blockrepeat{répéter \ovalnum{\textbf{4}} fois}
{
\blockmoreblocks{triangle équilatéral}
\blockmove{avancer de \ovalnum{50} pas}
\blockmove{\turnleft{} de \ovalnum{\textbf{90}} degrés}
}
\end{scratch}
&
\begin{scratch}[scale=0.9]
\blockinit{quand \greenflag est cliqué}
\blockmoreblocks{Initialisation}
\blockmoreblocks{triangle équilatéral}
\blockrepeat{répéter \ovalnum{\textbf{3}} fois}
{
\blockmove{avancer de \ovalnum{50} pas}
\blockmove{tourner \turnright{} \ovalnum{\textbf{60}} degrés}
\blockmoreblocks{triangle équilatéral}
\blockmove{tourner \turnright{} \ovalnum{\textbf{60}} degrés}
}
\end{scratch}\\ \hline
\end{tabularx}

\vspace{14pt}

\begin{tabularx}{\linewidth}{|*{3}{>{\centering \arraybackslash}X|}}\hline
\textbf{Figure A} & \textbf{Figure B} &\textbf{Figure C}\\ \hline
\begin{center}
%%début Tikz
%\begin{tikzpicture}[scale=0.4]
%% Triangle du bas
%\draw[thick] (0,0) -- (2,0) -- (1,1.732) -- cycle;
%% Triangle gauche
%\draw[thick] (0,0) -- (-2,0) -- (-1,1.732) -- cycle;
%% Triangle droit
%\draw[thick] (2,0) -- (4,0) -- (3,1.732) -- cycle;
%% Triangle du bas central inversé (ligne commune)
%% Petit triangle central pointe en bas
%\draw[thick] (0,1.732) -- (2,1.732) -- (1,0) -- cycle;
%\end{tikzpicture}
%Fin Tikz
%PSTricks
\psset{unit=1cm}
\begin{pspicture}(-1.5,-1.5)(1.5,1.5)
\pspolygon(1;-30)(1;90)(1;210)
\pspolygon(0.5;30)(0.5;150)(0.5;270)
\end{pspicture}
%fin PSTricks
\end{center}
\vspace{4pt}
&
\vspace{6pt}
\begin{center}
%%début Tikz
%\begin{tikzpicture}[scale=0.4]
%% Figure B : carré + triangles sur chaque côté, répété 4 fois (style losange)
%\draw[thick] (0,0) -- (2,0) -- (2,2) -- (0,2) -- cycle;
%\draw[thick] (0,0) -- (-1,-1.732) -- (1,-1.732) -- (2,0);
%\draw[thick] (2,0) -- (3.732,1) -- (2.732,2.732) -- (2,2);
%\end{tikzpicture}
%Fin Tikz
%PSTricks
\psset{unit=1cm}
\begin{pspicture}(-1.5,-1.5)(1.5,1.5)
%\pspolygon(1;-30)(1;90)(1;210)
\pspolygon(0.45;30)(0.45;150)(0.45;270)
\def\car{\psframe(-0.42,-0.42)(0.42,0.4)}
\multido{\n=90+120}{3}{\rput{\n}(0.66;\n){\car}}
\end{pspicture}
%fin PSTricks
\end{center}
\vspace{4pt}
&
\vspace{6pt}
\begin{center}
%%début Tikz
%\begin{tikzpicture}[scale=0.4]
%% Figure C : deux triangles se croisant (étoile)
%\draw[thick] (-1.5,0) -- (1.5,0) -- (0,2.6) -- cycle;
%\draw[thick] (-1.5,1.3) -- (1.5,1.3) -- (0,-1.3) -- cycle;
%\end{tikzpicture}
%Fin Tikz
%PSTricks
\psset{unit=1cm}
\begin{pspicture}(-1.5,-1.5)(1.5,1.5)
%\pspolygon(1;-30)(1;90)(1;210)
\def\tri{\pspolygon(0.57;-30)(0.57;90)(0.57;210)}
%\def\car{\psframe(-0.42,-0.42)(0.42,0.4)}
\psframe(-0.49,-0.49)(0.49,0.49)
%\multido{\n=0+90}{4}{\rput{\n}(0.66;\n){\tri}}
\rput(0,0.768){\tri}\rput{90}(-0.768,0){\tri}
\rput{180}(0,-0.768){\tri}\rput{-90}(0.768,0){\tri}
\end{pspicture}
%fin PSTricks
\end{center}
\vspace{6pt}\\ \hline
\end{tabularx}
\end{center}
\end{document}